\section{Limitations and Threats to Validity}
\label{sec:limitations}

The protocol effects reported here are large, and the temptation is to
read them as general. Several constraints bound that reading, stated
here before the conclusions.

\subsection{Scope of the Data}
\label{sec:limitations-scope}

All results come from three arrays at a single site. The arrays are
co-located and share one weather station, so their forecast errors are
correlated and consistency across them is consistency across module
technology under common meteorology, not geographic generalisation.
Every array-by-horizon table in this paper therefore contains fewer
independent observations than its cell count suggests, and no
significance test in this paper treats the nine cells as nine
independent trials.

The site is a hot desert climate with approximately 60 percent of
daylight hours classified as clear; results at a site with a different
cloud regime may differ, particularly the sky-stratified analysis in
Section \ref{sec:rq3}, where the partly cloudy class is the one in which
models add least. One evaluation year is used, at horizons of 1 to 6
hours, with deterministic forecasts only; Yang et al.
\cite{yang2020verification} argue that probabilistic verification
should be standard practice, and its absence here is a limitation
rather than a design choice.

The overcast sky class comprises 714 of 11,055 daylight hours,
approximately 6.5 percent; skill estimates for that class rest on
roughly one tenth the sample of the clear class, no significance test
was performed within sky-condition subsets, and the apparent ordering
of models under overcast conditions is not established.

\subsection{Bounds on the Architecture Results}
\label{sec:limitations-architecture-bounds}

Both architecture results in this paper are directional rather than
established, not demonstrated effects: full significance testing for
the convolutional comparison and the six-hour recurrent advantage is
reported in Section \ref{sec:rq2} and is not repeated here.

\subsection{The Residual Stage}
\label{sec:limitations-residual-stage}

The residual correction stage's measured cost is sensitive to a
construction detail - the two-fold-versus-four-fold sensitivity
reported in Section \ref{sec:rq2} - and that sensitivity is itself the
more robust finding.

Two further bounds apply. The fold-starvation mechanism was diagnosed
only for the recurrent base on sites 11 and 12; the convolutional base
shows a significant penalty in all nine cells including the one-hour
horizon, which that mechanism does not explain. And the five-year,
four-fold sensitivity run was never subjected to a Diebold-Mariano test:
its reported standard errors are two-sample errors across three seeds,
which measure training stochasticity rather than sampling uncertainty in
the evaluation period.

\subsection{Threats to Validity}
\label{sec:threats-to-validity}

\textit{Optimisation directive.} All models here are trained under a
squared error loss and evaluated by root-mean-square error and skill
scores derived from it. This is internally consistent, but Mayer
\cite{mayer2022physmlhybrid} demonstrates that the choice of directive
is itself a protocol axis: different error metrics are minimised by
different functionals, and running the same study under a
mean-absolute-error directive can reverse conclusions about whether a
component helps. We did not vary the directive, and results here should
not be assumed to transfer to a mean-absolute-error objective.

\textit{Weather regime.} The oracle regime is a perfect-forecast upper
bound, never presented as an achievable result; operational forecasting
sits between the lagged and oracle regimes, using numerical weather
prediction whose own error is represented in neither, so the gap
reported here bounds, rather than estimates, the value of weather
information in practice.

\textit{Reproducibility of recurrent results.} PyTorch provides no
deterministic CUDA backward pass for its recurrent layers (Section V),
which is why recurrent results are reported with seed variance across
five seeds rather than as point values.

\textit{Hyperparameters.} No hyperparameter tuning was performed, on any
split, for any model. This removes one class of leakage and makes the
architecture comparison a comparison of default configurations at
comparable scale, but it also means no model is presented at its best
achievable accuracy, and a differently tuned configuration might
reorder the architecture results. The protocol effects reported in
Section \ref{sec:rq1} are an order of magnitude larger than the
architecture differences and are unlikely to be affected.

\textit{Deferred work.} A second location and a probabilistic evaluation
were both scoped and deferred; the low-elevation clear-sky bias
motivating the daylight filter (Section \ref{sec:eval-protocol}) is a
candidate protocol axis not varied here. These are recorded in the
repository rather than claimed as future work in the abstract.
